\section{Background}

LGFVs emerged as organizational vehicles through which local governments could
finance development projects while navigating formal budget constraints. After
the 1994 tax sharing reform, local governments retained extensive expenditure
responsibilities but faced tighter formal revenue constraints. Platform
companies, land finance, and local state-owned enterprises became key
instruments for sustaining infrastructure investment and urban expansion
\citep{oi1992,montinola1995,jin2005,wong2013}.

The institutional logic of the platform model was straightforward. A local
government could keep formal budgetary borrowing limited while using a legally
separate company to raise funds, hold assets, develop land, and implement public
projects. The company could borrow against land-related revenue expectations,
project cash flows, asset injections, fiscal subsidies, or implicit public
support. This arrangement made local development finance more flexible, but it
also made the boundary between corporate debt and public obligation difficult
to observe. The same organizational form could contain commercial activity,
public infrastructure construction, fiscal prepayment, project repurchase, and
debt rollover.

The model expanded rapidly after the 2008 global financial crisis. Local
governments relied on LGFVs to borrow, build roads, develop industrial parks,
finance municipal utilities, and support land-centered urbanization. During
periods of rising land values and rapid growth, this model appeared viable.
Once the real estate market weakened and land transfer revenues contracted,
however, the same model exposed severe fiscal and financial fragility
\citep{tsui2011,bai2016longshadow}.

The central government responded by trying to move local borrowing from
platform-based channels into more formal fiscal instruments. The revised Budget
Law and State Council Document No. 43 established a framework in which local
governments could borrow through approved government bonds while unauthorized
borrowing and guarantees through platform companies were to be restricted
\citep{budgetlaw2014,statecouncil2014}. The policy shift did not simply ban
local debt. It attempted to distinguish legitimate government debt from
enterprise debt and to separate public fiscal responsibility from platform
company liabilities. This distinction became the institutional foundation for
later exit language: platform companies were expected to withdraw from
government financing functions, while local governments were expected to use
regulated debt instruments and budgetary channels.

Subsequent regulations tightened the boundary between local governments and
platform companies. The 2017 rules on local government financing behavior and
government purchase of services targeted guarantees, disguised borrowing,
irregular public-private arrangements, and the use of service-purchase contracts
as financing vehicles \citep{mof2017_50,mof2017_87}. These policies made the
formal rule increasingly clear: local governments should not use platform
companies to create new hidden debt, and platform companies should not rely on
government credit to perform financing functions. The practical question was
harder. Localities still had to finish projects, service existing obligations,
and sustain public investment under uneven fiscal conditions.

The current LGFV exit reform can therefore be understood as a central effort to
formalize local fiscal behavior rather than simply as a corporate restructuring
campaign. The central government seeks to bring local borrowing, investment, and
debt management back within more transparent fiscal and regulatory frameworks.
But local governments still face strong pressure to maintain investment, service
debt, and deliver development projects. This tension creates room for divergent
local responses. In some places, former platform functions can be absorbed by
special-purpose bonds, budgetary expenditure, regulated state assets, or
commercial operations. In others, the official platform label may change while
the underlying financing and project functions continue inside the same company
or move to another local state-owned entity.
